\documentclass[11pt,a4paper]{article}

% ── Packages ──
\usepackage[utf8]{inputenc}
\usepackage[T1]{fontenc}
\usepackage[english]{babel}
\usepackage[a4paper,margin=2.5cm]{geometry}
\usepackage{graphicx}
\usepackage{hyperref}
\usepackage{xcolor}
\usepackage{listings}
\usepackage{longtable}
\usepackage{array}
\usepackage{booktabs}
\usepackage{enumitem}
\usepackage{tabularx}
\usepackage{fancyhdr}
\usepackage{float}
\usepackage{parskip}


\hypersetup{
    colorlinks=true,
    linkcolor=blue!60!black,
    urlcolor=blue!70!black,
    citecolor=black,
    pdftitle={EmpPulse -- Detailed System Design},
    pdfauthor={Group 6}
}

% ── Code listings ──
\definecolor{codebg}{rgb}{0.96,0.96,0.96}
\definecolor{codegreen}{rgb}{0,0.5,0}
\definecolor{codeblue}{rgb}{0.15,0.15,0.65}
\definecolor{codepurple}{rgb}{0.55,0,0.75}

\lstset{
    basicstyle=\ttfamily\footnotesize,
    backgroundcolor=\color{codebg},
    breaklines=true,
    frame=single,
    showstringspaces=false,
    keywordstyle=\color{codeblue}\bfseries,
    stringstyle=\color{codepurple},
    commentstyle=\color{codegreen}\itshape,
    tabsize=2,
    xleftmargin=0.3cm,
    aboveskip=0.8em,
    belowskip=0.8em
}

\lstdefinestyle{json}{
    basicstyle=\ttfamily\footnotesize,
    backgroundcolor=\color{codebg},
    breaklines=true,
    frame=single,
    showstringspaces=false,
    stringstyle=\color{codepurple},
    morestring=[b]",
    morekeywords={true,false,null}
}

\lstdefinestyle{mermaid}{
    basicstyle=\ttfamily\footnotesize,
    backgroundcolor=\color{codebg},
    breaklines=true,
    frame=single,
    showstringspaces=false
}

% ── Headers ──
\pagestyle{fancy}
\fancyhf{}
\fancyhead[L]{\small EmpPulse}
\fancyhead[R]{\small Detailed System Design}
\fancyfoot[C]{\thepage}

% ══════════════════════════════════════════════════════════════
\begin{document}

% ── Title page ──
\begin{titlepage}
\centering
\vspace*{3cm}
{\Huge\bfseries EmpPulse\par}
\vspace{0.5cm}
{\LARGE Detailed System Design\par}
\vspace{2cm}
{\Large Group 6\par}
\vspace{1cm}
{\large\today\par}
\vfill
\end{titlepage}

\tableofcontents
\newpage

% ══════════════════════════════════════════════════════════════
\section{Introduction}

EmpPulse is an employee management system that enables organizations to track leave requests, log working hours, and manage departments. The system supports three roles---Employee, Administrator, and Owner---each with distinct permissions. This document provides the detailed system design, expanding upon the initial system design with precise architectural decisions, module internals, data models, API contracts, UI flows, security mechanisms, error handling, and deployment configuration.

% ══════════════════════════════════════════════════════════════
\section{Refined System Architecture}
\label{sec:architecture}

\subsection{Overview}

EmpPulse follows a monolithic Spring Boot architecture with a React-based frontend served as static assets. The backend exposes a REST API consumed by the single-page application (SPA). Session-based authentication (JSESSIONID) secures all protected endpoints. An external SMTP service handles transactional emails (credential delivery, password resets, notification emails). Both the Java application and PostgreSQL database run inside Docker containers on a single server.

\subsection{Architecture Diagram}

\begin{figure}[H]
    \centering
    \includegraphics[width=0.8\linewidth]{architectureSS.png}
    \label{fig:placeholder}
\end{figure}

\subsection{Component Responsibilities}

\begin{itemize}[leftmargin=1.5em]
    \item \textbf{React Frontend} -- Single-page application served as static assets; communicates with the backend exclusively through REST API calls with session cookies.
    \item \textbf{REST API Layer} -- Spring MVC controllers; receives HTTP requests, delegates to service layer, returns JSON responses.
    \item \textbf{Authentication Module} -- Spring Security configuration using HTTP session authentication (\texttt{JSESSIONID} cookie). Handles login, logout, session validation, and password reset token management.
    \item \textbf{User Management Module} -- CRUD for Users, Employees, Admins; enforces role-based creation rules and visibility constraints.
    \item \textbf{Department Management Module} -- CRUD for Departments; manages Admin--Department and Employee--Department associations.
    \item \textbf{Leave Request Module} -- Full lifecycle of leave requests including creation, modification, modification-requests, approval/rejection, cancellation, and vacation balance updates.
    \item \textbf{Logged Hours Module} -- Creation, editing, deletion of work hour logs; implements overlap/adjacency merging logic.
    \item \textbf{Default Hours Module} -- Manages default weekly hour templates for Departments and Employees.
    \item \textbf{Export Module} -- Generates CSV data archives scoped to the requesting Admin's departments.
    \item \textbf{Notification Module} -- Sends transactional emails via SMTP (credentials, password changes, leave status updates).
    \item \textbf{Scheduler} -- A \texttt{@Scheduled} Spring component that runs at midnight to auto-generate logged hours for the following day based on default hours configuration.
    \item \textbf{Validation Layer} -- Shared validators enforcing business rules (date ranges, email format, uniqueness, ownership checks).
    \item \textbf{PostgreSQL Database} -- Persistent relational storage running in a Docker container with a mounted volume for data durability.
\end{itemize}

% ══════════════════════════════════════════════════════════════
\newpage
\section{Detailed Module Design}
\label{sec:modules}

The backend is organized package-by-feature under \texttt{com.oman.EmpPulse}. Each module is split into three sub-packages: \texttt{web} (Spring MVC \texttt{@RestController} classes), \texttt{internal} (services, JPA entities, and \texttt{Repository} interfaces), and \texttt{dto} (request/response payloads). Cross-module communication does not reach into another module's \texttt{internal} package; instead each module publishes a small \texttt{api} package containing an interface (e.g. \texttt{UserApi}, \texttt{LeaveApi}, \texttt{NotificationApi}) plus the DTOs it shares. The service classes implement these \texttt{*Api} interfaces. Request flow is therefore \texttt{Controller} $\rightarrow$ \texttt{Service (\textit{*Api} impl)} $\rightarrow$ \texttt{Repository}. Validation is performed inline inside the services (there are no standalone \texttt{Validator} classes); the method signatures listed below are taken directly from the source.

\subsection{Module Overview Diagram}

\begin{figure}[H]
    \centering
    \includegraphics[width=0.5\linewidth]{modularSS.png}
    \caption{Enter Caption}
    \label{fig:placeholder}
\end{figure}

% ── 3.1 Authentication Module ──
\subsection{Authentication Module}

\textbf{Responsibility:} Handles user sign-in, sign-out, session management, and password reset flows.

\textbf{Internal components:}
\begin{itemize}[leftmargin=1.5em]
    \item \texttt{auth.web.AuthController} -- Endpoints under \texttt{/api/auth}: \texttt{login}, \texttt{logout}, \texttt{forgotPassword}, \texttt{resetPassword}.
    \item \texttt{auth.internal.AuthService} -- \texttt{authenticate(email, password)} verifies credentials against the User module via \texttt{UserApi} and the configured \texttt{PasswordEncoder}.
    \item \texttt{auth.internal.PasswordResetService} -- \texttt{requestReset(email)}, \texttt{validateAndConsume(rawToken)}, \texttt{reset(rawToken, newPassword)}.
    \item \texttt{auth.internal.PasswordResetToken} / \texttt{PasswordResetTokenRepository} -- Persist time-limited reset tokens (15-minute expiry).
\end{itemize}

\textbf{Data owned:} \texttt{Password\_reset\_token} table. The session itself is the Spring Security context persisted server-side; \texttt{AuthController.login} builds a \texttt{UsernamePasswordAuthenticationToken} and saves it via the \texttt{SecurityContextRepository}, which issues the \texttt{JSESSIONID} cookie.

\textbf{Public operations (controller endpoints):}
\begin{itemize}[leftmargin=1.5em]
    \item \texttt{login(LoginRequest)} -- Authenticates; on success establishes the security context and returns the caller's profile, otherwise \texttt{401} with code \texttt{INVALID\_CREDENTIALS}.
    \item \texttt{logout()} -- Invalidates the current session.
    \item \texttt{forgotPassword(PasswordForgotRequest)} -- Delegates to \texttt{PasswordResetService.requestReset}; always returns \texttt{200}.
    \item \texttt{resetPassword(PasswordResetRequest)} -- Delegates to \texttt{PasswordResetService.reset}.
\end{itemize}

\textbf{Dependencies:} User Management Module (\texttt{UserApi} for credential lookup and password reset), Notification Module (\texttt{NotificationApi} for reset emails).

\textbf{Business rules:}
\begin{itemize}[leftmargin=1.5em]
    \item Invalid credentials return \texttt{401 Unauthorized} without specifying which field was wrong.
    \item Reset token expires after 15 minutes; new request invalidates previous token.
    \item Password change invalidates all other active sessions for that user.
\end{itemize}

\textbf{Error cases:} Invalid credentials, expired/invalid reset token, email not found.

% ── 3.2 User Management Module ──
\subsection{User Management Module}

\textbf{Responsibility:} CRUD operations for Users, Employees, Admins, and Bonus Vacation Days. Enforces role-based creation, visibility, and modification rules.

\textbf{Internal components:}
\begin{itemize}[leftmargin=1.5em]
    \item \texttt{user.web.UserController} -- Profile and user CRUD endpoints (\texttt{/api/me}, \texttt{/api/users}).
    \item \texttt{user.web.AdminController} -- \texttt{GET /api/admins} (\texttt{listAdmins}).
    \item \texttt{user.web.EmployeeController} -- \texttt{GET /api/employees} (\texttt{listEmployees}).
    \item \texttt{user.internal.UserService} (\texttt{implements UserApi}) -- Orchestrates User + Employee/Admin creation, profile loading, password and preference changes.
    \item \texttt{user.internal.EmployeeService} (\texttt{implements EmployeeApi}) -- Employee lookups, admin-access checks, week-schedule assignment.
    \item \texttt{user.internal.AdminService} (\texttt{implements AdminApi}) -- Admin lookups, department-oversight queries, owner resolution.
    \item \texttt{UserRepository}, \texttt{EmployeeRepository}, \texttt{AdminRepository}, \texttt{BonusVacationDaysRepository}; entities \texttt{User}, \texttt{Employee}, \texttt{Admin}, \texttt{BonusVacationDays}.
\end{itemize}

\textbf{Data owned/modified:} \texttt{User}, \texttt{Employee}, \texttt{Admin}, \texttt{Bonus\_vacation\_days} tables.

\textbf{Public operations:}
\begin{itemize}[leftmargin=1.5em]
    \item \textit{UserController:} \texttt{getMe}, \texttt{changePassword}, \texttt{updatePreferences}, \texttt{createUser}, \texttt{getUserProfile}, \texttt{updateUser}, \texttt{deleteUser}, \texttt{getBonusVacationDays}, \texttt{updateBonusVacationDays}.
    \item \textit{UserService:} \texttt{createUser(req, callerUserId, callerIsOwner)}, \texttt{getUserProfile(userId, callerId, callerIsOwner)}, \texttt{updateUser(...)}, \texttt{softDeleteUser(userId)}, \texttt{changeMyPassword(userId, req, currentSessionId)}, \texttt{updatePreferences(userId, req)}, \texttt{resetPassword(userId, rawNewPassword)}, \texttt{updateBonusVacationDays(userId, req, callerUserId)}; plus \texttt{UserApi} methods \texttt{findActiveByEmail}, \texttt{loadProfile}, \texttt{findContactById}, \texttt{ensureOwnerExists}.
    \item \textit{EmployeeService:} \texttt{getEmployeesForAdmin(id)}, \texttt{requireAdminAccessToEmployee(adminId, employeeId)}, \texttt{findActiveEmployeeIds()}, \texttt{hasEmployeesInDepartment(departmentId)}, \texttt{getWeekScheduleId} / \texttt{assignWeekSchedule}.
    \item \textit{AdminService:} \texttt{getAllAdmins()}, \texttt{getOwnerAdmin()}, \texttt{departmentIdsForAdminUser(adminId)}, \texttt{overseesDepartment(adminId, departmentId)}.
\end{itemize}

\textbf{Dependencies:} Department Module, Notification Module, Leave Request Module (cascading deactivation). Soft-deletion (\texttt{softDeleteUser}) detaches the user and removes dependent data rather than hard-deleting.

\textbf{Business rules:}
\begin{itemize}[leftmargin=1.5em]
    \item User and Employee/Admin must be created in a single atomic operation.
    \item Active users must have unique email among active users.
    \item Only Owner can create Admins; Admins can create Employees.
    \item Owner can only be a User who is Admin for all Departments.
    \item Users, Employees, and Admins are never hard-deleted (soft deactivation only).
    \item Deactivating an Employee removes them from their department and deletes pending leave requests.
    \item Deactivating an Admin detaches them from all departments.
\end{itemize}

% ── 3.3 Department Management Module ──
\subsection{Department Management Module}

\textbf{Responsibility:} CRUD for Departments, management of Admin--Department and Employee--Department associations.

\textbf{Internal components:}
\begin{itemize}[leftmargin=1.5em]
    \item \texttt{department.web.DepartmentController} -- Department CRUD endpoints under \texttt{/api/departments}.
    \item \texttt{department.internal.DepartmentService} (\texttt{implements DepartmentApi}) -- Business logic for department operations and admin--department associations.
    \item \texttt{DepartmentRepository}; entity \texttt{Department} (admin links handled inside \texttt{DepartmentService.setAdminDepartments}).
\end{itemize}

\textbf{Data owned/modified:} \texttt{Department} table and the admin--department association it manages.

\textbf{Public operations:}
\begin{itemize}[leftmargin=1.5em]
    \item \textit{DepartmentController:} \texttt{listDepartments}, \texttt{createDepartment}, \texttt{getDepartment}, \texttt{updateDepartment}, \texttt{deleteDepartment}.
    \item \textit{DepartmentService:} \texttt{getDepartmentsForAdmin(adminId)}, \texttt{getDepartment(departmentId)}, \texttt{createDepartment(req)}, \texttt{updateDepartment(departmentId, req)}, \texttt{deleteDepartment(departmentId)}, \texttt{setAdminDepartments(adminId, departmentIds)}, \texttt{requireAdminAccessToDepartment(adminId, departmentId)}, \texttt{overseesDepartment}/\texttt{isAdminOfDepartment}, \texttt{ensureDefaultDepartmentExists()}, \texttt{getWeekScheduleId} / \texttt{assignWeekSchedule}.
\end{itemize}

\textbf{Business rules:}
\begin{itemize}[leftmargin=1.5em]
    \item Department names must be unique.
    \item Default Department cannot be deleted.
    \item Owner can detach Admin only if that Admin oversees more than 2 departments.
    \item Admin can move Employees only across departments they oversee.
    \item Admin cannot attach/detach themselves or other Admins.
\end{itemize}

% ── 3.4 Leave Request Module ──
\subsection{Leave Request Module}

\textbf{Responsibility:} Full leave request lifecycle: creation, modification, modification requests, approval, rejection, cancellation, deletion. Manages vacation balance updates.

\textbf{Internal components:}
\begin{itemize}[leftmargin=1.5em]
    \item \texttt{leave.web.LeaveController} -- Leave request CRUD, response, modification, and cancel endpoints under \texttt{/api/leave-requests}.
    \item \texttt{leave.internal.LeaveService} (\texttt{implements LeaveApi}) -- Core business logic for leave operations, including inline date-range/overlap validation and vacation-balance computation (\texttt{countUsedVacationDays}). There are no separate \texttt{LeaveValidator} or \texttt{VacationBalanceService} classes.
    \item \texttt{LeaveRepository}; entity \texttt{Leave} with \texttt{LeaveType} and \texttt{LeaveStatus} enums.
\end{itemize}

\textbf{Data owned/modified:} \texttt{Leave} table; reads \texttt{Employee} and \texttt{Bonus\_vacation\_days} for balance calculation.

\textbf{Public operations:}
\begin{itemize}[leftmargin=1.5em]
    \item \textit{LeaveController:} \texttt{listLeaveRequests}, \texttt{getLeaveRequest}, \texttt{createLeaveRequest}, \texttt{updateLeaveRequest}, \texttt{respondToLeaveRequest} (\texttt{PATCH .../response}, carries approve/reject), \texttt{createModification} (\texttt{POST .../modifications}), \texttt{cancelLeaveRequest} (\texttt{PATCH .../cancel}), \texttt{deleteLeaveRequest}.
    \item \textit{LeaveService:} \texttt{createLeaveRequest(req, callerId, isAdmin)}, \texttt{listLeaveRequests(callerId, isAdmin)}, \texttt{getLeaveRequest(id, callerId)}, \texttt{updateLeaveRequest(...)}, \texttt{respondToLeaveRequest(...)}, \texttt{createModification(...)}, \texttt{cancelLeaveRequest(id, callerId)}, \texttt{deleteLeaveRequest(id, callerId)}; plus \texttt{LeaveApi} methods \texttt{findActiveLeavesByEmployeeIds}, \texttt{countUsedVacationDays(employeeId, year)}, \texttt{findEmployeeIdsOnUnpaidApprovedLeave}.
\end{itemize}

\textit{Note: approval and rejection share a single \texttt{respondToLeaveRequest} operation rather than separate \texttt{approve}/\texttt{reject} endpoints.}

\textbf{Business rules:}
\begin{itemize}[leftmargin=1.5em]
    \item Start date must not be after end date.
    \item Active leave requests must not overlap with other active requests for the same Employee (unless linked as modification).
    \item Personal leave type requires a reason.
    \item Default paid type: Vacation and Sick are paid; Personal is unpaid.
    \item Employee-created requests start as Pending; if Employee is also Admin, auto-Accepted.
    \item Admin-created requests are auto-Accepted.
    \item Only Employee can cancel their own Accepted request.
    \item Only Employee can hard-delete their own Pending request.
    \item Modification request to an Accepted leave creates a linked Pending leave.
    \item Accepted modification overwrites original and self-deletes.
    \item Rejected modification unlinks and becomes standalone Rejected request.
    \item Accepting an unpaid leave deletes all logged hours for that Employee in the leave date range.
    \item Vacation balance is updated upon creation, cancellation, rejection, and date changes of Vacation-type leaves.
\end{itemize}

% ── 3.5 Logged Hours Module ──
\subsection{Logged Hours Module}

\textbf{Responsibility:} Creation, editing, and deletion of work hour logs, including overlap/adjacency merge logic.

\textbf{Internal components:}
\begin{itemize}[leftmargin=1.5em]
    \item \texttt{loggedhours.web.LoggedHoursController} -- CRUD endpoints under \texttt{/api/employees/\{employeeId\}/logged-hours}.
    \item \texttt{loggedhours.internal.LoggedHoursService} (\texttt{implements LoggedHoursApi}) -- Business logic, inline validation (\texttt{ensureRequiredFieldsPresent}, \texttt{validateTimeRange}, \texttt{validateNotFuture}) and overlap/adjacency merge logic; no separate \texttt{MergeService} class.
    \item \texttt{LoggedHoursRepository}; entity \texttt{LoggedHours}.
\end{itemize}

\textbf{Data owned/modified:} \texttt{Logged\_Hours} table.

\textbf{Public operations:}
\begin{itemize}[leftmargin=1.5em]
    \item \textit{LoggedHoursController:} \texttt{listLoggedHours}, \texttt{createLoggedHours}, \texttt{editLoggedHours}, \texttt{deleteLoggedHours}.
    \item \textit{LoggedHoursService:} \texttt{listLoggedHours(employeeId, callerId, isAdmin)}, \texttt{createLoggedHours(...)} (triggers merge), \texttt{updateLoggedHours(...)} (triggers merge), \texttt{deleteLoggedHours(employeeId, loggedHoursId, callerId)}; plus \texttt{LoggedHoursApi} methods \texttt{deleteByEmployeeAndDateRange} and \texttt{createAutoLogIfAbsent} (used by the scheduler).
\end{itemize}

\textbf{Business rules:}
\begin{itemize}[leftmargin=1.5em]
    \item End time must be after start time.
    \item Future dates are not allowed.
    \item Employees cannot create, modify, or delete logged hours---only Admins can.
    \item On creation or edit, if the new interval overlaps or is adjacent to existing entries for the same Employee on the same date, all affected entries are merged into a single entry using the earliest start and latest end time.
\end{itemize}

% ── 3.6 Default Hours Module ──
\subsection{Default Hours Module}

\textbf{Responsibility:} Manages default weekly hour templates for Departments and Employees. Provides data for the automatic hour logging scheduler.

\textbf{Internal components:}
\begin{itemize}[leftmargin=1.5em]
    \item \texttt{defaulthours.web.DefaultHoursController} -- GET/PUT \texttt{/api/employees/\{employeeId\}/default-hours}.
    \item \texttt{defaulthours.web.DepartmentDefaultHoursController} -- GET/PUT \texttt{/api/departments/\{departmentId\}/default-hours}.
    \item \texttt{defaulthours.internal.DefaultHoursService} (\texttt{implements DefaultHoursApi}) -- Resolves effective default hours (Employee-level overrides Department-level).
    \item \texttt{WeekScheduleRepository}, \texttt{ScheduleBlockRepository}; entities \texttt{WeekSchedule} and \texttt{ScheduleBlock}.
\end{itemize}

\textbf{Data owned/modified:} \texttt{Week\_schedule} and \texttt{Schedule\_block} tables (a \texttt{WeekSchedule} groups the per-day \texttt{ScheduleBlock} intervals).

\textbf{Public operations:}
\begin{itemize}[leftmargin=1.5em]
    \item \textit{Controllers:} \texttt{getDefaultHours} / \texttt{setDefaultHours} for both employee and department scopes.
    \item \textit{DefaultHoursService:} \texttt{getEmployeeDefaultHours(employeeId, callerAdminId)}, \texttt{setEmployeeDefaultHours(...)}, \texttt{getDepartmentDefaultHours(departmentId, callerAdminId)}, \texttt{setDepartmentDefaultHours(...)}, \texttt{inheritDepartmentSchedule(departmentId)}, \texttt{findEmployeeIntervalForDay(...)} (consumed by the scheduler).
\end{itemize}

\textbf{Business rules:}
\begin{itemize}[leftmargin=1.5em]
    \item Each day of week can have multiple time intervals.
    \item End time must be after start time for each interval.
    \item Days default to N/A (no intervals) until configured.
    \item If Employee has no personal default hours, they inherit the Department's defaults.
    \item If Department has no default hours, Employee has no inherited default.
\end{itemize}

% ── 3.7 Scheduler Module ──
\subsection{Scheduler Module}

\textbf{Responsibility:} Automatically generates logged hours at midnight for the following day.

\textbf{Internal components:}
\begin{itemize}[leftmargin=1.5em]
    \item \texttt{loggedhours.internal.AutoLogScheduler} -- A Spring \texttt{@Scheduled} component (it resides in the Logged Hours module). Entry point \texttt{scheduledAutoLog()} is bound to a configurable cron expression and zone (\texttt{emppulse.autolog.cron} / \texttt{emppulse.autolog.zone}) and delegates to \texttt{autoLogForToday()}.
\end{itemize}

\textbf{Operation:} \texttt{autoLogForToday()} iterates active employees (\texttt{EmployeeApi.findActiveEmployeeIds}), resolves each employee's interval for the day (\texttt{DefaultHoursApi.findEmployeeIntervalForDay}), skips employees on approved unpaid leave (\texttt{LeaveApi.findEmployeeIdsOnUnpaidApprovedLeave}), and calls \texttt{LoggedHoursApi.createAutoLogIfAbsent} for each interval.

\textbf{Business rules:}
\begin{itemize}[leftmargin=1.5em]
    \item Creates one \texttt{Logged\_Hours} entry per configured interval.
    \item If the day has no configured intervals, nothing is created.
    \item Does not create entries for Employees with approved unpaid leave on that day.
    \item Existing entries are not duplicated (the \texttt{createAutoLogIfAbsent} guard).
\end{itemize}

% ── 3.8 Export Module ──
\subsection{Export Module}

\textbf{Responsibility:} Generates downloadable CSV data archives.

\textbf{Implementation status:} \textit{Not yet implemented.} No \texttt{export} package, controller, or service currently exists in the backend, and no \texttt{/api/exports} endpoint is registered. The module is retained here as planned scope; the design below describes the intended structure.

\textbf{Planned components:}
\begin{itemize}[leftmargin=1.5em]
    \item \texttt{ExportController} -- \texttt{GET /api/exports/csv} endpoint.
    \item \texttt{ExportService} -- Queries and formats data into CSV, archives into ZIP.
\end{itemize}

\textbf{Business rules:}
\begin{itemize}[leftmargin=1.5em]
    \item Only Administrators can export.
    \item Exports are scoped to the departments the Admin oversees.
    \item Data may include users, leaves, logged hours, and department information.
    \item Output format is \texttt{.csv} files archived into a \texttt{.zip}.
\end{itemize}

% ── 3.9 Notification Module ──
\subsection{Notification Module}

\textbf{Responsibility:} Sends transactional emails via external SMTP.

\textbf{Internal components:}
\begin{itemize}[leftmargin=1.5em]
    \item \texttt{notification.internal.NotificationService} (\texttt{implements NotificationApi}) -- Entry point other modules call; wraps each send in \texttt{sendSafely(...)} so an SMTP failure does not abort the primary operation.
    \item \texttt{notification.internal.EmailMessageFactory} -- Builds the \texttt{EmailMessage(subject, body)} for each trigger (using \texttt{app.base-url} for links).
    \item \texttt{notification.internal.EmailSender} -- \texttt{send(to, subject, body)} wrapper over Spring \texttt{JavaMailSender}.
\end{itemize}

\textbf{Public operations (\texttt{NotificationApi}):} \texttt{sendAccountCredentials}, \texttt{sendPasswordResetEmail}, \texttt{sendPasswordChangedNotification}, \texttt{sendEmailChangedNotification}, \texttt{sendLeaveCreatedOnBehalf}, \texttt{sendLeaveDecision}, \texttt{sendLeaveModifiedByAdmin}.

\textbf{Email triggers (mapped to operations):}
\begin{itemize}[leftmargin=1.5em]
    \item Account credentials for a newly created user -- \texttt{sendAccountCredentials}.
    \item Password reset link -- \texttt{sendPasswordResetEmail}.
    \item Password change confirmation -- \texttt{sendPasswordChangedNotification}.
    \item Email/credentials changed by Owner -- \texttt{sendEmailChangedNotification}.
    \item Leave request created on behalf of Employee -- \texttt{sendLeaveCreatedOnBehalf}.
    \item Leave request approved/rejected -- \texttt{sendLeaveDecision}.
    \item Leave request modified by Admin -- \texttt{sendLeaveModifiedByAdmin}.
\end{itemize}

% ══════════════════════════════════════════════════════════════
\section{Detailed Data Model / Database Design}
\label{sec:datamodel}

The database schema is implemented in PostgreSQL. The full schema is defined in \texttt{db\_design.txt};

\subsection{Enumerations}

\begin{longtable}{p{3.5cm} p{10cm}}
\toprule
\textbf{Enum} & \textbf{Values} \\
\midrule
\texttt{theme} & \texttt{DARK}, \texttt{LIGHT} \\
\texttt{language} & \texttt{ENG}, \texttt{UKR} \\
\texttt{leave\_type} & \texttt{SICK}, \texttt{VACATION}, \texttt{PERSONAL} \\
\texttt{leave\_status} & \texttt{PENDING}, \texttt{REJECTED}, \texttt{APPROVED}, \texttt{CANCELED} \\
\bottomrule
\end{longtable}

\subsection{Data model}

\begin{figure}[H]
    \centering
    \includegraphics[width=1\linewidth]{Data model.png}
    \caption{Graphical representation of system data mode}
    \label{fig:data_model}
\end{figure}

\section{API Design}
\label{sec:api}

All endpoints are prefixed with \texttt{/api}. Authentication is required for all endpoints except the auth endpoints. The session cookie \texttt{JSESSIONID} is automatically sent by the browser. Responses use JSON format.

\subsection{Authentication}
\begin{description}
    \item[\texttt{POST /api/auth/login}] Sign in and create a session.
    \item[\texttt{POST /api/auth/logout}] Sign out and invalidate the current session.
    \item[\texttt{POST /api/auth/password/forgot}] Request a password reset email.
    \item[\texttt{POST /api/auth/password/reset}] Reset password with token.
\end{description}

\subsection{Users}
\begin{description}
    \item[\texttt{GET /api/me}] Get my profile.
    \item[\texttt{POST /api/me/password}] Change password (logged in).
    \item[\texttt{PATCH /api/me/preferences}] Update my preferences.
    \item[\texttt{POST /api/users}] Create a user (\textit{owner only}).
    \item[\texttt{GET /api/users/\{userId\}}] Get user details.
    \item[\texttt{PATCH /api/users/\{userId\}}] Update user.
    \item[\texttt{DELETE /api/users/\{userId\}}] Soft delete user (\textit{owner only}).
\end{description}

\subsection{Departments}
\begin{description}
    \item[\texttt{GET /api/departments}] List departments.
    \item[\texttt{POST /api/departments}] Create department (\textit{owner only}).
    \item[\texttt{GET /api/departments/\{departmentId\}}] Get department details.
    \item[\texttt{PATCH /api/departments/\{departmentId\}}] Update department.
    \item[\texttt{DELETE /api/departments/\{departmentId\}}] Delete department (\textit{owner only}).
    \item[\texttt{GET /api/departments/\{departmentId\}/default-hours}] Get department default week hours.
    \item[\texttt{PUT /api/departments/\{departmentId\}/default-hours}] Set department default week hours.
\end{description}

\subsection{Employees}
\begin{description}
    \item[\texttt{GET /api/employees}] List employees.
    \item[\texttt{GET /api/employees/\{employeeId\}/default-hours}] Get employee default week hours.
    \item[\texttt{PUT /api/employees/\{employeeId\}/default-hours}] Set employee default week hours.
\end{description}

\subsection{Leave Requests}
\begin{description}
    \item[\texttt{GET /api/leave-requests}] List leave requests.
    \item[\texttt{POST /api/leave-requests}] Create leave request.
    \item[\texttt{GET /api/leave-requests/\{leaveRequestId\}}] Get leave request details.
    \item[\texttt{PATCH /api/leave-requests/\{leaveRequestId\}}] Update a leave request.
    \item[\texttt{DELETE /api/leave-requests/\{leaveRequestId\}}] Delete PENDING leave request (\textit{employee only}).
    \item[\texttt{POST /api/leave-requests/\{leaveRequestId\}/cancel}] Cancel a leave request (\textit{employee only}).
    \item[\texttt{POST /api/leave-requests/\{leaveRequestId\}/approve}] Approve a leave request.
    \item[\texttt{POST /api/leave-requests/\{leaveRequestId\}/reject}] Reject a leave request.
\end{description}

\subsection{Logged Hours}
\begin{description}
    \item[\texttt{GET /api/employees/\{employeeId\}/logged-hours}] List logged work hours.
    \item[\texttt{POST /api/employees/\{employeeId\}/logged-hours}] Create logged hours (\textit{admin only}).
    \item[\texttt{PATCH /api/employees/\{employeeId\}/logged-hours/\{loggedHoursId\}}] Update logged hours (\textit{admin only}).
    \item[\texttt{DELETE /api/employees/\{employeeId\}/logged-hours/\{loggedHoursId\}}] Delete logged hours (\textit{admin only}).
\end{description}

\subsection{Exports}
\begin{description}
    \item[\texttt{GET /api/exports/csv}] Export CSV data.
\end{description}

% ══════════════════════════════════════════════════════════════
\section{User Interface Design}
\label{sec:ui}

The UI is designed in Figma. This section describes the main screens, their purpose, user actions, and the API endpoints they consume. Full visual mockups are available in the linked \href{https://www.figma.com/design/gXW9SjcVd3WYGhZkZgdiW4/CRM-Dashboard-Customers-List--Community-?node-id=0-1&t=B8TP7vt4psiaqP2p-1}{Figma project}

% ══════════════════════════════════════════════════════════════
\section{Detailed Sequence and Activity Diagrams}
\label{sec:flows}

\subsection{Flow 1: User Login}

\begin{figure}[H]
    \centering
    \includegraphics[width=1\linewidth]{flow1.png}
    \label{fig:placeholder}
\end{figure}

\subsection{Flow 2: Employee Creates Leave Request}

\begin{figure}[H]
    \centering
    \includegraphics[width=1\linewidth]{flow2.png}
    \label{fig:placeholder}
\end{figure}

\subsection{Flow 3: Admin Approves / Rejects Leave Request}

\begin{figure}[H]
    \centering
    \includegraphics[width=1\linewidth]{flow3.png}
    \label{fig:placeholder}
\end{figure}

\subsection{Flow 4: Logged Hours Creation with Merge}

\begin{figure}[H]
    \centering
    \includegraphics[width=1\linewidth]{flow4.png}
    \label{fig:placeholder}
\end{figure}

\subsection{Flow 5: Leave Request State Diagram}

\begin{figure}[H]
    \centering
    \includegraphics[width=1\linewidth]{flow6.png}
    \label{fig:placeholder}
\end{figure}

\newpage

\subsection{Flow 6: Automatic Hour Logging (Scheduler)}

\begin{figure}[!htbp]
    \centering
    \includegraphics[width=\linewidth,height=0.85\textheight,keepaspectratio]{diagram5.png}
    \label{fig:placeholder}
\end{figure}

\newpage
% ══════════════════════════════════════════════════════════════
\section{Business Logic and Validation Rules}
\label{sec:logic}

The full detailed rules are defined in \texttt{business\_logic.md};

% ══════════════════════════════════════════════════════════════
\section{Security Design}
\label{sec:security}

\subsection{Authentication}

\begin{itemize}[leftmargin=1.5em]
    \item \textbf{Method:} HTTP session-based authentication using Spring Security.
    \item \textbf{Session identifier:} \texttt{JSESSIONID} cookie.
    \item \textbf{Password storage:} Bcrypt hash. Plaintext passwords are never stored.
    \item \textbf{Session lifecycle:}
    \begin{itemize}
        \item Created upon successful \texttt{POST /api/auth/login}.
        \item Invalidated upon \texttt{POST /api/auth/logout}.
        \item All sessions invalidated upon password change.
        \item Server-side session timeout configurable (default: 30 minutes of inactivity).
    \end{itemize}
    \item \textbf{Password reset:}
    \begin{itemize}
        \item Token-based via email; token expires after 15 minutes.
        \item New reset request invalidates previous token.
        \item After reset, all existing sessions are invalidated.
    \end{itemize}
\end{itemize}

\subsection{Authorization and Roles}

The system has three roles with hierarchical permissions:

\begin{longtable}{p{2.5cm} p{11cm}}
\toprule
\textbf{Role} & \textbf{Permissions} \\
\midrule
\textbf{Employee} &
    View own profile, leave requests, and logged hours.
    Create and manage own leave requests (create, edit pending, cancel accepted, delete pending).
    Change own password.
    Update UI preferences (theme, language). \\
\midrule
\textbf{Admin} &
    All Employee permissions for themselves.
    View Employees and their data in overseen Departments.
    Create Employee accounts in overseen Departments.
    Create, edit, approve, reject leave requests for overseen Employees.
    Create, edit, delete logged hours for overseen Employees.
    Manage default hours for overseen Departments and their Employees.
    Move Employees across overseen Departments.
    Export CSV data for overseen Departments.
    Manage Bonus Vacation Days for overseen Employees. \\
\midrule
\textbf{Owner} &
    All Admin permissions across all Departments.
    Create Admin accounts.
    Create, modify, delete Departments.
    Attach/detach Admins to/from Departments.
    Modify all User/Employee/Admin fields.
    Deactivate Employees and Admins.
    View all Admins and Users. \\
\bottomrule
\end{longtable}

\subsection{Access Control Implementation}

\begin{itemize}[leftmargin=1.5em]
    \item Every protected endpoint requires a valid \texttt{JSESSIONID}.
    \item Unauthenticated requests to protected endpoints receive \texttt{401 Unauthorized}.
    \item The backend resolves the current user from the session and checks role + department oversight before processing any request.
    \item Resource ownership is validated server-side: an Admin can only access data for Employees in Departments they oversee. Violations return \texttt{403 Forbidden}.
    \item The frontend hides UI elements based on role, but this is a UX convenience only---the backend enforces all rules independently.
\end{itemize}

\subsection{Input Validation and Protection}

\begin{itemize}[leftmargin=1.5em]
    \item All input is validated on both frontend (immediate feedback) and backend (authoritative).
    \item Email format validated using standard regex patterns.
    \item Date/time fields validated for logical consistency (start before end, no future dates for logged hours).
    \item SQL injection prevented by using JPA/Hibernate parameterized queries.
    \item XSS mitigated by React's default escaping and backend input sanitization.
    \item CSRF protection via Spring Security's built-in CSRF token mechanism (synchronizer token pattern), coordinated with the session cookie.
    \item Login error messages are generic (``Invalid credentials'') to prevent email enumeration.
    \item Password reset always returns 200 OK regardless of whether the email exists.
\end{itemize}

% ══════════════════════════════════════════════════════════════
\section{Error Handling and Edge Cases}
\label{sec:errors}

\subsection{Error Response Format}

The backend produces two distinct error shapes depending on where the error originates.

\textbf{1. Authentication and access-control errors.} The login, password-forgot and
password-reset endpoints (\texttt{AuthController}), together with the global
authentication-entry-point and access-denied handlers (\texttt{SecurityConfig}), return a
compact JSON object carrying a machine-readable \texttt{code} and a human-readable
\texttt{message}:

\begin{lstlisting}[style=json]
{
  "code": "INVALID_CREDENTIALS",
  "message": "Invalid credentials"
}
\end{lstlisting}

\textbf{2. Domain and business-rule errors.} Every service signals these by throwing
\texttt{ResponseStatusException(status, reason)}, which Spring Boot serializes through its
default error handler at \texttt{/error}. Because \texttt{server.error.include-message} is
left at its default (\texttt{never}), the \texttt{reason} is logged and attached to the
exception but is \emph{not} written to the response body; the body carries only the standard
error attributes:

\begin{lstlisting}[style=json]
{
  "timestamp": "2026-05-12T10:30:00.000+00:00",
  "status": 409,
  "error": "Conflict",
  "path": "/api/leave-requests"
}
\end{lstlisting}

The React frontend therefore does not depend on the server-supplied message for these errors:
it maps the HTTP status code to a localized, user-facing message (see \texttt{services/api.ts}),
applying per-endpoint overrides where a more specific wording is warranted.

\textit{Note: input validation is performed imperatively inside the services (there is no
Jakarta Bean Validation / \texttt{@Valid} on the controllers), so there is no multi-field
\texttt{fieldErrors} response -- each invalid input yields a single \texttt{400 Bad Request}
with the corresponding reason listed below.}

\subsection{Error Catalog}

The catalog below lists the error cases the implementation actually raises. For the
authentication block the \textbf{Response} column gives the \texttt{code} and \texttt{message}
returned in the JSON body; for every other block it gives the \texttt{reason} attached to the
\texttt{ResponseStatusException} (carried by the exception, surfaced to the user by the
frontend via the status code).

\begin{longtable}{p{6cm} p{1.5cm} p{6cm}}
\toprule
\textbf{Situation} & \textbf{HTTP Code} & \textbf{Response (code / reason)} \\
\midrule
\endhead
\multicolumn{3}{l}{\textit{Authentication and Password Reset (code/message JSON)}} \\
\midrule
Invalid email or password on login & 401 & \texttt{INVALID\_CREDENTIALS} -- ``Invalid credentials'' (no hint about which field is wrong). \\
Unauthenticated request to a protected endpoint & 401 & \texttt{UNAUTHORIZED} -- ``Unauthorized''. Frontend treats this as an expired session and redirects to login. \\
Forgot password with a blank email & 400 & \texttt{EMAIL\_REQUIRED} -- ``Email is required''. \\
Forgot password for an unknown account & 404 & \texttt{USER\_NOT\_FOUND} -- ``No account found for that email''. \\
Reset password with blank password/confirmation & 400 & \texttt{PASSWORD\_REQUIRED} -- ``New password and confirmation are required''. \\
Reset password confirmation mismatch & 400 & \texttt{PASSWORD\_MISMATCH} -- ``Passwords do not match''. \\
Reset token invalid, expired or already used & 400 & \texttt{INVALID\_TOKEN} -- ``Reset link is invalid or expired''. \\
\midrule
\multicolumn{3}{l}{\textit{Authorization Errors}} \\
\midrule
Authenticated user lacks the role required by an endpoint & 403 & \texttt{FORBIDDEN} -- ``Forbidden'' (Spring Security access-denied handler). \\
Admin accesses an employee / department / leave request outside their departments & 403 & ``No access to this employee'' / ``No access to this department'' / ``No access to this leave request''. \\
Admin tries to create a non-employee account & 403 & ``Admins can only create employee accounts''. \\
Admin edits owner-only user fields & 403 & ``Admins can only update employee department and vacation balance''. \\
Non-owner detaches an employee from a department & 403 & ``Only the owner can detach an employee from a department''. \\
Acting on a leave request the caller does not own (delete / cancel / modify) & 403 & ``Only the employee the request belongs to can delete / cancel / modify it''. \\
Admin targets a department they do not oversee & 403 & ``Admin does not oversee target department''. \\
Generic access fallback & 403 & ``Access denied'' / ``Not an admin''. \\
\midrule
\multicolumn{3}{l}{\textit{Validation Errors (400 Bad Request)}} \\
\midrule
Create-user payload missing required fields & 400 & ``Name, surname, email and password are required''. \\
Create user with no role & 400 & ``Cannot create an account with no role''. \\
Invalid preference value & 400 & ``Invalid theme'' / ``Invalid language''. \\
Change-password payload incomplete & 400 & ``Current password, new password and confirmation are required''. \\
Current password wrong / new equals old / confirmation mismatch (self-service change) & 400 & ``Current password is incorrect'' / ``New password must differ from the current password'' / ``Passwords do not match''. \\
Negative vacation balance / bonus days & 400 & ``Yearly vacation balance must be greater or equal to 0'' / ``Bonus vacation days must be greater or equal to 0''; ``Year and days are required''. \\
Leave-create payload missing required fields & 400 & ``employeeId, type, paid, startDate and endDate are required''. \\
Leave end date before start date & 400 & ``End date must not be before start date''. \\
Reason missing for a personal leave & 400 & ``Reason is required for personal leave''. \\
Non-admin supplies an admin comment & 400 & ``adminComment is allowed only for administrators''. \\
Invalid decision on a leave response & 400 & ``Decision must be approved or rejected''. \\
Modification changes nothing & 400 & ``A modification must change at least one field''. \\
Logged-hours payload incomplete / inverted / in the future & 400 & ``date, startTime and endTime are required'' / ``Start time must be before end time'' / ``Cannot log hours for a future date''. \\
Default-hours payload invalid & 400 & ``days is required''; ``dayOfWeek must be between 0 and 6''; ``Duplicate dayOfWeek: N''; ``intervals is required''; ``Only one interval per day is allowed''; ``startTime and endTime are required''; ``Start time must be before end time''. \\
\midrule
\multicolumn{3}{l}{\textit{Conflict / Business Rule Violations (409 Conflict)}} \\
\midrule
Email already used by another account & 409 & ``Email already in use''. \\
Attempt to delete the owner & 409 & ``Owner cannot be deleted''. \\
Leave dates overlap an existing request & 409 & ``Overlapping leave request exists''. \\
Delete a non-pending leave request & 409 & ``Only pending leave requests can be deleted''. \\
Approve/reject a non-pending request & 409 & ``Only pending leave requests can be approved or rejected''. \\
Cancel a non-approved request & 409 & ``Only approved leave requests can be cancelled''. \\
Propose a modification on a non-approved request & 409 & ``Modifications can only be proposed for approved requests''. \\
Second modification on a request that already has one pending & 409 & ``A pending modification already exists for this request'' / ``Request has a pending modification; resolve it first''. \\
Edit a rejected/cancelled request & 409 & ``Rejected or cancelled requests cannot be edited''. \\
Edit an approved request in place & 409 & ``Approved requests are edited via a modification request''. \\
File leave for an inactive employee & 409 & ``Employee is not active''. \\
Duplicate department name & 409 & ``Department name already in use''. \\
Delete the default department & 409 & ``Cannot delete the default department''. \\
Delete a department with employees/admins assigned & 409 & ``Cannot delete department: employees are assigned'' / ``Cannot delete department: administrators are assigned''. \\
Reassignment would leave an admin with no department & 409 & ``Cannot remove admin N from their only department here; to deactivate this admin, edit their account instead''. \\
\midrule
\multicolumn{3}{l}{\textit{Resource Not Found (404)}} \\
\midrule
Referenced entity does not exist or is out of the caller's scope & 404 & ``User not found'' / ``Employee not found'' / ``Department not found'' / ``Leave request not found'' / ``Logged hours entry not found'' / ``Admin not found''. \\
\midrule
\multicolumn{3}{l}{\textit{Internal Server Errors (500)}} \\
\midrule
A required entity referenced by an invariant cannot be resolved (e.g. owner, an authenticated user's department) & 500 & ``Data inconsistency''. Signals a broken invariant rather than user error; logged server-side. \\
Unhandled failure (database connectivity, mail transport, etc.) & 500 & Falls through to Spring's default error response; logged server-side. Email sending is best-effort -- a failed send is logged and does not fail the primary operation. \\
\bottomrule
\end{longtable}

% ══════════════════════════════════════════════════════════════
\section{Deployment and Runtime View}
\label{sec:deployment}

\subsection{Overview}

The EmpPulse application is deployed on a single server using Docker Compose to orchestrate two containers: one for the Java (Spring Boot) application and one for PostgreSQL. The React frontend is built into static assets and served by the Spring Boot application (or optionally by a reverse proxy). External connectivity is provided over HTTPS on port 443.

\subsection{Deployment Architecture}

The deployment architecture is shown in the attached diagram (Figure~\ref{fig:deployment}) and described below:

\begin{itemize}[leftmargin=1.5em]
    \item \textbf{Server:} A single Linux server hosts the entire application stack.
    \item \textbf{Docker Compose:} Orchestrates the container environment.
    \item \textbf{Java Container:} Runs the Spring Boot application (backend + static frontend assets). Exposes port 443 (HTTPS) to the Internet. Connects to PostgreSQL via JDBC on the internal Docker network (port 5432).
    \item \textbf{PostgreSQL Container:} Runs PostgreSQL database. Not exposed to the Internet; accessible only from the Java container on the internal Docker network.
    \item \textbf{DB Volume:} A Docker named volume mounted to the PostgreSQL container ensures data persistence across container restarts and redeployments.
    \item \textbf{External SMTP:} The Java container connects outbound to an external SMTP server for sending transactional emails.
\end{itemize}

\begin{figure}[H]
    \centering
    \includegraphics[width=0.7\linewidth]{deployment_diagram.png}
    \caption{Deployment Diagram}
    \label{fig:deployment_diagram}
\end{figure}

\subsection{Environment Configuration}

Configuration is provided through environment variables, injected via Docker Compose \texttt{.env} file or environment section:

\begin{longtable}{p{5.5cm} p{8cm}}
\toprule
\textbf{Variable} & \textbf{Purpose} \\
\midrule
\texttt{DB\_HOST} & PostgreSQL hostname (default: \texttt{postgres}). \\
\texttt{DB\_PORT} & PostgreSQL port (default: \texttt{5432}). \\
\texttt{DB\_NAME} & Database name (e.g., \texttt{emppulse}). \\
\texttt{DB\_USER} & Database username. \\
\texttt{DB\_PASSWORD} & Database password. \\
\texttt{SMTP\_HOST} & SMTP server hostname. \\
\texttt{SMTP\_PORT} & SMTP server port. \\
\texttt{SMTP\_USER} & SMTP authentication username. \\
\texttt{SMTP\_PASSWORD} & SMTP authentication password. \\
\texttt{APP\_BASE\_URL} & Public URL of the application (for email links). \\
\texttt{SESSION\_TIMEOUT} & HTTP session timeout (default: 30m). \\
\bottomrule
\end{longtable}

\subsection{Technology Stack}

\begin{longtable}{p{3.5cm} p{10cm}}
\toprule
\textbf{Component} & \textbf{Technology} \\
\midrule
Frontend & React (Vite build tooling), served as static assets. \\
Backend & Java, Spring Boot (Spring MVC, Spring Security, Spring Data JPA, Spring Mail, Spring Scheduler). \\
Database & PostgreSQL (latest stable), running in Docker. \\
ORM & Hibernate / Spring Data JPA. \\
Authentication & Spring Security with HTTP session (JSESSIONID). \\
Email & Spring Mail (\texttt{JavaMailSender}) with external SMTP. \\
Containerization & Docker, Docker Compose. \\
Build & Maven (backend), npm/Vite (frontend). \\
\bottomrule
\end{longtable}


\end{document}